\section{Self-orthogonal embeddings}\label{ch:so}

\subsection{Preliminaries}

Self-orthogonal codes have been studied for their combinatorial properties and applications in lattice theory~\cite{Harada2009, Bouyuklieva2013}, $t$-designs~\cite{Hadara2004, Bachoc2004}, and quantum error correction~\cite{Calderbank1998, Kim2002:2}. Pless~\cite{Pless1968} enumerated self-orthogonal codes and used finite geometry to prove the uniqueness of the Golay codes. Subsequent work has focused on classifying optimal self-orthogonal codes across various fields and rings~\cite{Pless1972, Huffman2005, Boukllieve2006}.

Recently, the problem of embedding a linear code into a self-orthogonal code has gained attention. Kim et al.~\cite{Kim2021} introduced the concept of \textit{self-orthogonal embedding} (SO embedding), devising algorithms to find the shortest self-orthogonal supercode for codes of dimension up to 4. This was extended by Kim and Choi~\cite{Kim2022} to dimensions 5 and 6 using self-orthogonality matrices. In this section, we generalize these results, leveraging the geometric properties of the \textit{hull} of a code to determine the exact length of shortest SO embeddings for arbitrary binary linear codes.

\begin{definition}[Self-orthogonal and self-dual codes]
The \textit{dual} of a linear code $\mathcal{C}$, denoted by $\mathcal{C}^\perp$, is the set of vectors orthogonal to every codeword in $\mathcal{C}$ with respect to the standard inner product. A linear code $\mathcal{C}$ is called \textit{self-orthogonal} if $\mathcal{C} \subseteq \mathcal{C}^\perp$, and \textit{self-dual} if $\mathcal{C} = \mathcal{C}^\perp$.
\end{definition}

\begin{definition}[Hull]
The \textit{hull} of a linear code $\mathcal{C}$, denoted by $\Hull(\mathcal{C})$, is the intersection of the code and its dual:
\begin{center}
$\Hull(\mathcal{C}) = \mathcal{C} \cap \mathcal{C}^\perp.$
\end{center}
The dimension of the hull is denoted by $\ell = \dim(\Hull(\mathcal{C}))$. The hull plays a crucial role in determining the complexity of embedding $\mathcal{C}$ into a self-orthogonal code.
\end{definition}

\begin{definition}[Self-orthogonal embedding]
Let $\mathcal{C}$ be a binary $[n,k]$ code. An \textit{SO embedding} of $\mathcal{C}$ is a self-orthogonal code $\tilde{\mathcal{C}}$ of length $n' \ge n$ such that puncturing $\tilde{\mathcal{C}}$ on a specific set of $n' - n$ coordinates yields $\mathcal{C}$. An SO embedding is called a \textit{shortest SO embedding} if its length $n'$ is minimal among all SO embeddings of $\mathcal{C}$.
\end{definition}

A fundamental property linking generator matrices to self-orthogonality is the Gram matrix condition.

\begin{proposition}[Gram matrix condition]
Let $G$ be a generator matrix of a linear code $\mathcal{C}$. The code $\mathcal{C}$ is self-orthogonal if and only if the Gram matrix $GG^T$ is the zero matrix over the base field, i.e., $GG^T = \mathbf{0}$.
\end{proposition}

\begin{definition}[Hermitian inner product]
For codes over $\mathbb{F}_{q^2}$, the Hermitian inner product of two vectors $x, y \in \mathbb{F}_{q^2}^n$ is defined as $\langle x, y \rangle_H = \sum_{i=1}^n x_i y_i^q$. A code is Hermitian self-orthogonal if $\langle c, c' \rangle_H = 0$ for all $c, c' \in \mathcal{C}$. The condition for the generator matrix becomes $G \bar{G}^T = \mathbf{0}$, where $\bar{G}$ denotes the element-wise $q$-th power (conjugation).
\end{definition}

\subsection{Rank-one elimination algorithm}
The discovery of self-orthogonal embeddings represents a significant theoretical contribution of the SolEvolve framework. The problem asks for the shortest self-orthogonal code that contains a given linear code as a subcode. The Generator conceived a novel algorithmic approach to this problem, deriving the Rank-one elimination strategy from first principles.

The core insight begins with the observation that a generator matrix $G$ defines a self-orthogonal code if and only if the Gram matrix $GG^T$ is the zero matrix over the base field. For an arbitrary linear code with generator matrix $G$, the product $GG^T$ is generally non-zero. The goal is to append columns to $G$ such that the new matrix $G' = [G \mid S]$ satisfies $G'(G')^T = \mathbf{0}$. This condition expands to $GG^T + SS^T = \mathbf{0}$, or equivalently, $SS^T = GG^T$ (since we are working in characteristic 2, addition and subtraction are identical).

The algorithm proceeds by iteratively eliminating the non-zero entries of the residual matrix $R = GG^T$. In each step, the Generator identifies a non-zero element in $R$. If a diagonal element $R_{ii} = 1$ is non-zero, it implies that the $i$-th row of $G$ is not self-orthogonal. To correct this, we must append a column that has a 1 in the $i$-th position. If an off-diagonal element $R_{ij}$ is non-zero, it implies that rows $i$ and $j$ are not orthogonal. The algorithm constructs a specific column vector $v$ to append to $S$ such that the new product $SS^T$ cancels out this specific non-zero entry in $R$. This process is repeated until $R$ becomes the zero matrix.

The algorithm achieves provable optimality. The number of added columns is exactly equal to the rank of the initial Gram matrix $GG^T$. Thus, the length of the shortest self-orthogonal embedding is $n + \rank(GG^T)$. This constructive proof not only solves the embedding problem but also provides an efficient $O(k^2n + rk^2)$ algorithm for constructing the embedding, which is significantly faster than brute-force search.

\begin{theorem}
Let $\mathcal{C}$ be a binary $[n, k]$ code with generator matrix $G$, and let $A = GG^T$ be its Gram matrix with $r = \rank(A)$.
\begin{enumerate}
    \item[(i)] Any SO embedding $\widetilde{G} = [G|S]$ requires $S$ to have at least $r$ columns.
    \item[(ii)] There always exists an $S \in \F_2^{k \times r}$ such that $\widetilde{G} = [G|S]$ is self-orthogonal.
\end{enumerate}
Therefore, the shortest length of a self-orthogonal embedding of $\mathcal{C}$ is $n+r$.
\end{theorem}

\begin{proof}
(i) The condition for self-orthogonality is $A = SS^T$. By basic properties of matrix rank, $\rank(SS^T) \le \rank(S)$. The rank of $S$ is at most its number of columns, say $s$. Thus, $r = \rank(A) = \rank(SS^T) \le s$. The number of columns $s$ must be at least $r$.

(ii) The existence of such an $S$ is guaranteed by the theory of quadratic forms over $\F_2$. The symmetric matrix $A$ corresponds to a quadratic form $q(x) = x^T A x$. By the Witt decomposition theorem, $A$ is congruent to a direct sum of simple blocks. An explicit construction of $S$ with $r$ columns is provided by the algorithm in the next subsection.
\end{proof}

\begin{algorithm}[tb]
\caption{\textsc{Rank-OneElimination}}
\label{alg:rank1}
\begin{algorithmic}
\REQUIRE A generator matrix $G \in \F_2^{k \times n}$.
\ENSURE An extended matrix $\widetilde{G}=[G\mid S]$ such that $\widetilde{G}\widetilde{G}^{T}=0$.
\STATE \textit{// Notation: $R_{*,i}$ = column $i$ of $R$; $e_i \in \F_2^k$ = standard basis vector (1 at position $i$, 0 elsewhere)}
\STATE \textit{// Initialize residual matrix $R$ as the Gram matrix; $S$ collects appended columns}
\STATE $R \gets GG^T$ \hfill \textit{// $R \in \F_2^{k \times k}$: symmetric residual to be eliminated}
\STATE $S \gets [\,]$ \hfill \textit{// $S$: column matrix to be appended to $G$}
    \WHILE{$R \neq \mathbf{0}$}
        \IF{$\exists i$ such that $R_{ii} = 1$}
            \STATE Find any such $i$. \hfill \textit{// Anisotropic case: diagonal entry is non-zero}
            \STATE $v \gets R_{*,i}$ \hfill \textit{// $v$: the $i$-th column of $R$, a $k$-dimensional vector}
        \ELSE
            \STATE Find any $(i,j)$ with $i<j$ and $R_{ij}=1$. \hfill \textit{// Hyperbolic case}
            \STATE $v \gets e_i$ \hfill \textit{// $e_i$: standard basis vector with 1 in position $i$}
        \ENDIF
        \STATE $S \gets [S \mid v]$ \hfill \textit{// Append $v$ as a new column to $S$}
        \STATE $R \gets R + vv^T$ \hfill \textit{// Rank-one update: eliminates targeted entry from $R$}
    \ENDWHILE
\STATE $\widetilde{G} \gets [G \mid S]$ \hfill \textit{// Final extended generator matrix}
\end{algorithmic}
\end{algorithm}

\paragraph{Variable Definitions for Algorithm~\ref{alg:rank1}.} To clarify the algorithm's operation: $R \in \F_2^{k \times k}$ is the \textit{residual matrix} initialized to the Gram matrix $GG^T$; the algorithm iteratively reduces $R$ to zero. The notation $R_{*,i}$ denotes the $i$-th column of $R$, and $e_i \in \F_2^k$ is the standard basis vector with a single 1 in position $i$. The matrix $S$ accumulates the columns appended to $G$, with each column $v$ chosen to perform a rank-one update $R \gets R + vv^T$ that eliminates either a diagonal entry (anisotropic case) or prepares for subsequent elimination (hyperbolic case). The update $R \gets R + vv^T$ exploits the characteristic-2 property where addition equals subtraction.

\begin{theorem}[Correctness and optimality]
\label{thm:main-exact-length}
Algorithm~\ref{alg:rank1} terminates after exactly $r = \rank(GG^T)$ iterations and returns a generator matrix $\widetilde{G}$ of a self-orthogonal code of the minimum possible length $n+r$.
\end{theorem}

\begin{proof}[Proof sketch]
The core idea is that each step correctly reduces the complexity of the residual matrix $R$.
\begin{itemize}
    \item When $R_{ii}=1$ (anisotropic line), choosing $v=R_{*,i}$ and updating $R \to R+vv^T$ can be shown to decrease the rank of the non-singular part of $R$ by one.
    \item When all $R_{ii}=0$ and $R_{ij}=1$ (hyperbolic pair), choosing $v=e_i$ updates $R \to R' = R + e_ie_i^T$. This changes the diagonal: $R'_{ii}$ becomes $1$. The next iteration will then fall into the first case, and it can be shown that these two combined steps reduce the rank by 2.
\end{itemize}
The total number of iterations corresponds to decomposing $A$ into its Witt components. Let $A$ have $a$ anisotropic lines and $h$ hyperbolic pairs, so $r = a+2h$. The algorithm performs $a$ single-step eliminations and $h$ two-step eliminations, for a total of $a+2h=r$ iterations. Thus, $S$ has exactly $r$ columns. By Theorem~\ref{thm:main-exact-length}, this is optimal. Upon termination, $R=\mathbf{0}$, which implies $GG^T + SS^T = \mathbf{0}$, so the resulting code is self-orthogonal.
\end{proof}

\subsection{SAT-based self-orthogonal embeddings}

This section consolidates the binary and non-binary SAT experiments on shortest self-orthogonal (SO) embeddings. While the Rank-one elimination algorithm provides a theoretical upper bound and a constructive method, it has a significant limitation: it cannot manage the minimum distance of the resulting code. For example, in the binary case, the Rank-one algorithm yielded an embedding with $d=4$. In contrast, our SolEvolve SAT strategy successfully incorporated minimum distance constraints, increasing the distance by 2 to achieve $d=6$ for the binary $[22,11]$ code. Furthermore, we demonstrated that this SAT-based approach is not limited to $\mathbb{F}_2$ but generalizes effectively to $\mathbb{F}_3, \mathbb{F}_4,$ and $\mathbb{F}_5$. The encodings use one-hot symbols, explicit field tables (addition, multiplication, Hermitian conjugation where applicable), and, when tractable, explicit minimum-distance constraints. We report both optimality with respect to the Gram-rank lower bound and proximity to Grassl’s distance tables.

Table~\ref{tab:so-results} integrates the binary Hamming studies with the ternary BCH, quaternary Hermitian, and quinary dot-product experiments. Each line lists the Gram-rank lower bound $t_{\min}$, the smallest $t$ solved by SAT, the attained minimum distance, and the Grassl bound $u_b=l_b$ when available.



Over the binary field $\mathbb{F}_2$, the system embedded the Hamming $[15,11,3]$ code into the shortened Golay $[22,11,6]$ code. This is a classic result, but SolEvolve rediscovered it autonomously.

We also investigated the possibility of achieving a minimum distance of $d=7$.

\begin{theorem}[Non-existence of self-orthogonal embedding with $d=7$]
The optimal self-orthogonal embedding of the Hamming code $[15,11,3]$ is strictly limited to $d=6$. Specifically, no self-orthogonal embedding exists with $d \ge 7$.
\end{theorem}

\begin{proof}
The SAT solver was configured to search for a self-orthogonal embedding of $\mathcal{H}_4$ with $d \ge 7$. This search involved 71,845 variables and 282,503 clauses. The solver returned \textbf{UNSATISFIABLE} after 300 seconds. This result provides a computer-assisted proof that the optimal self-orthogonal embedding of the Hamming code $[15,11,3]$ is strictly limited to $d=6$, even though linear codes with parameters $[22,11,7]$ are known to exist (e.g., by shortening the Golay code). This highlights the structural constraint imposed by the self-orthogonality requirement.
\end{proof}

More interestingly, over the ternary field $\mathbb{F}_3$, it extended a BCH $[13,7,5]$ code to a binary $[20,7,9]$ code self-orthogonal code. This is the first reported instance of this specific construction being found via a SAT-based approach.

\begin{proposition}[Optimality of Ternary BCH Embedding]
We analyze the optimality of this embedding following the rank-one elimination framework. The theoretical minimum number of added columns $t$ is lower-bounded by the rank of the Gram matrix $GG^T$. For the ternary BCH $[13,7,5]$ code, we computed $\rank(GG^T) = 7$ over $\mathbb{F}_3$. Thus, any self-orthogonal embedding must add at least $t \ge 7$ columns. Our SAT solver successfully found a solution at exactly $t=7$, achieving the shortest possible embedding length $n' = n + t = 13 + 7 = 20$.

The discovered code has minimum Hamming distance $d=9$ and its weight enumerator polynomial is:
\begin{align*}
W(z) =\; & 1 + 166z^9 + 972z^{12} \\
&+ 954z^{15} + 94z^{18}.
\end{align*}
With $t=7$, the extended generator $G_{\text{ext}}$ is:
\begin{center}
\resizebox{\linewidth}{!}{%
$G_{\text{ext}} = \left[ \begin{array}{ccccccccccccc|ccccccc}
1 & 0 & 2 & 0 & 2 & 2 & 1 & 0 & 0 & 0 & 0 & 0 & 0 & 2 & 0 & 0 & 0 & 1 & 1 & 1 \\
0 & 1 & 0 & 2 & 0 & 2 & 2 & 1 & 0 & 0 & 0 & 0 & 0 & 0 & 2 & 0 & 0 & 2 & 2 & 0 \\
0 & 0 & 1 & 0 & 2 & 0 & 2 & 2 & 1 & 0 & 0 & 0 & 0 & 0 & 0 & 2 & 0 & 2 & 1 & 0 \\
0 & 0 & 0 & 1 & 0 & 2 & 0 & 2 & 2 & 1 & 0 & 0 & 0 & 0 & 0 & 0 & 2 & 1 & 2 & 0 \\
0 & 0 & 0 & 0 & 1 & 0 & 2 & 0 & 2 & 2 & 1 & 0 & 0 & 0 & 0 & 0 & 0 & 0 & 1 & 1 \\
0 & 0 & 0 & 0 & 0 & 1 & 0 & 2 & 0 & 2 & 2 & 1 & 0 & 0 & 0 & 0 & 0 & 2 & 2 & 1 \\
0 & 0 & 0 & 0 & 0 & 0 & 1 & 0 & 2 & 0 & 2 & 2 & 1 & 0 & 0 & 0 & 0 & 1 & 2 & 2
\end{array} \right]$%
}
\end{center}
\end{proposition}

The framework also handled extension fields. For the Hermitian code over $\mathbb{F}_4$, it constructed a binary $[15,5,6]$ code self-orthogonal embedding starting from a random $[10,5]$ code. The resulting generator matrix is:
\begin{center}
\resizebox{\linewidth}{!}{%
$G_{\text{ext}} = \left[ \begin{array}{cccccccccc|ccccc}
0 & \omega^2 & \omega & 1 & 1 & \omega^2 & 0 & \omega & 0 & 0 & 0 & 1 & 1 & 1 & \omega^2 \\
\omega & \omega^2 & \omega & \omega^2 & \omega & \omega^2 & \omega & 0 & \omega^2 & 1 & 1 & 0 & 1 & 0 & \omega^2 \\
\omega & 1 & 0 & \omega^2 & \omega^2 & \omega & 1 & \omega^2 & \omega & 1 & \omega^2 & \omega & 0 & 0 & \omega^2 \\
1 & 0 & 0 & \omega & \omega^2 & 0 & \omega^2 & \omega^2 & 1 & \omega & 1 & \omega & 0 & 0 & \omega^2 \\
0 & \omega^2 & \omega & 1 & 0 & \omega^2 & 1 & \omega^2 & \omega & \omega^2 & \omega & 0 & 0 & 0 & 1
\end{array} \right]$%
}
\end{center}
The weight enumerator polynomial is:
\begin{align*}
W(z) =\; & 1 + 12z^6 + 69z^8 + 372z^{10} \\
&+ 402z^{12} + 168z^{14}.
\end{align*}
Note that we use the notation $\{0, 1, \omega, \omega^2\}$ for elements of $\mathbb{F}_4$, where $\omega^2 = \omega + 1$.

Similarly, over $\mathbb{F}_5$, it found an optimal binary $[18,6,8]$ code embedding. The current best at $t=6$ uses:
\begin{center}
\resizebox{\linewidth}{!}{%
$G_{\text{ext}} = \left[ \begin{array}{cccccccccccc|cccccc}
0 & 3 & 2 & 0 & 4 & 1 & 1 & 0 & 1 & 0 & 1 & 4 & 1 & 2 & 0 & 0 & 0 & 1 \\
2 & 4 & 2 & 1 & 3 & 4 & 4 & 4 & 0 & 2 & 1 & 1 & 1 & 4 & 2 & 1 & 0 & 0 \\
1 & 4 & 3 & 1 & 2 & 3 & 0 & 3 & 2 & 4 & 3 & 1 & 0 & 0 & 4 & 0 & 4 & 2 \\
4 & 3 & 1 & 2 & 3 & 1 & 1 & 0 & 0 & 2 & 0 & 2 & 1 & 3 & 4 & 2 & 1 & 0 \\
2 & 0 & 1 & 0 & 0 & 2 & 0 & 3 & 1 & 4 & 0 & 3 & 3 & 4 & 3 & 3 & 2 & 2 \\
2 & 4 & 3 & 4 & 2 & 1 & 0 & 4 & 1 & 3 & 2 & 1 & 0 & 4 & 3 & 3 & 4 & 3
\end{array} \right]$%
}
\end{center}
with weight enumerator polynomial:
\begin{align*}
W(z) =\; & 1 + 4z^8 + 48z^9 + 232z^{10} \\
&+ 436z^{11} + 1436z^{12} + 2152z^{13} \\
&+ 3752z^{14} + 3172z^{15} + 2824z^{16} \\
&+ 1304z^{17} + 264z^{18}.
\end{align*}
and $d_{\min}=8$. These results highlight the complementary strengths of the two approaches. While the Rank-one elimination algorithm provides a constructive proof for the minimum length, it is limited by its inability to manage the minimum distance (yielding only $d=4$ for the binary case). In contrast, the SolEvolve SAT strategy successfully incorporated minimum distance constraints, increasing the distance by 2 to achieve $d=6$ in $\mathbb{F}_2$, and demonstrated robust generalization across $\mathbb{F}_3, \mathbb{F}_4,$ and $\mathbb{F}_5$. 

